% ======================================================================
%  PART 5 -- Trends, Seasonality and Differencing: SARIMA
% ======================================================================
\section{Trends, Seasonality and Differencing: SARIMA}
\label{sec:sarima}

Real signals rarely look stationary: they drift (a \emph{trend}) and they
repeat (a \emph{season}). The working model of this chapter is
$X_t=\mu_t+Y_t$ with $\mu_t$ a deterministic time-dependent mean and $\{Y_t\}$
stationary. There are two ways to peel off $\mu_t$: \emph{estimate and subtract}
(regression on $t$ and on seasonal dummies), or \emph{difference it away} with
the operator $\diff=I-B$. Differencing of order $d$ annihilates a polynomial
trend of degree $d-1$; a \emph{seasonal} difference $\diff_{(s)}=I-B^{s}$ kills
a period-$s$ season. Layering autoregressive and moving-average structure onto
\emph{both} the ordinary lags and the seasonal lags gives the SARMA family, and
folding the differencing in front gives the (S)ARIMA model
$\Phi(B)(1-B)^{d}X_t=\Theta(B)\eps_t$ --- the standard Box--Jenkins workhorse.
Throughout, $\{\eps_t\}$ is mean-zero white noise of variance $\sigma_\eps^2$,
AR/seasonal-AR coefficients are $\phi_j,\phi^{(s)}_j$ and MA/seasonal-MA
coefficients are $\theta_k,\theta^{(s)}_k$.

% ----------------------------------------------------------------------
\subsection{Definitions}

\begin{definition}{Mean-trend (signal-plus-noise) model}{p5-trendmodel}
A series with a slowly varying level is modelled as
\[
X_t=\mu_t+Y_t,
\]
where $\mu_t$ is a deterministic, time-dependent mean (the \textbf{trend}) and
$\{Y_t\}$ is a zero-mean second-order stationary process. The canonical linear
case is $\mu_t=a+bt$; more generally $\mu_t$ may be a degree-$(q-1)$ polynomial
in $t$. When the level instead repeats, $\mu_t=s_t$ with $s_{t+s}=s_t$ is a
\textbf{periodic} (seasonal) trend of period $s$; the two can coexist,
$X_t=a+bt+s_t+Y_t$.
\end{definition}

\begin{definition}{Difference operator $\diff$}{p5-diff}
The (first-)\textbf{difference operator} is
\[
\diff=I-B,\qquad \diff X_t=X_t-X_{t-1},
\]
where $B$ is the backshift operator $BX_t=X_{t-1}$. It is \emph{linear}. Its
$d$-th power is, by the binomial theorem,
\[
\diff^{d}=(I-B)^{d}=\sum_{k=0}^{d}\binom{d}{k}(-1)^{k}B^{k},
\qquad
\diff^{d}X_t=\sum_{k=0}^{d}\binom{d}{k}(-1)^{k}X_{t-k}.
\]
For example $\diff^{2}X_t=X_t-2X_{t-1}+X_{t-2}$.
\end{definition}

\begin{definition}{Seasonal difference operator $\diff_{(s)}$}{p5-seasdiff}
The \textbf{seasonal difference} of period $s$ is
\[
\diff_{(s)}=I-B^{s},\qquad \diff_{(s)}X_t=X_t-X_{t-s}.
\]
It removes a deterministic component of period $s$ (since
$\diff_{(s)}s_t=s_t-s_{t-s}=0$) and, applied to a linear trend, leaves a
constant: $\diff_{(s)}(bt)=bs$. The ordinary difference is the case $s=1$,
$\diff=\diff_{(1)}$. To kill a linear-plus-seasonal trend one composes them,
$\diff\diff_{(s)}=(I-B)(I-B^{s})$.
\end{definition}

\begin{definition}{Seasonal moving average $\mathrm{SMA}(Q)_s$}{p5-sma}
A \textbf{seasonal moving average} of order $Q$ at period $s$ is
\[
X_t=\eps_t-\sum_{j=1}^{Q}\theta^{(s)}_j\,\eps_{t-js}
=\Theta^{(s)}(B)\eps_t,
\qquad
\Theta^{(s)}(z)=1-\sum_{j=1}^{Q}\theta^{(s)}_j\,z^{sj}.
\]
Only lags that are \emph{multiples of $s$} appear. The basic seasonal MA(1) is
$X_t=\eps_t-\theta\eps_{t-s}$.
\end{definition}

\begin{definition}{Seasonal autoregression $\mathrm{SAR}(P)_s$}{p5-sar}
A \textbf{seasonal autoregression} of order $P$ at period $s$ is
\[
X_t=\eps_t+\sum_{j=1}^{P}\phi^{(s)}_j\,X_{t-js},
\qquad
\Phi^{(s)}(z)=1-\sum_{j=1}^{P}\phi^{(s)}_j\,z^{sj},
\qquad
\Phi^{(s)}(B)X_t=\eps_t .
\]
Again only the seasonal lags $s,2s,\dots,Ps$ enter.
\end{definition}

\begin{definition}{SARMA$(p,q)\times(P,Q)_s$}{p5-sarma}
Combining an ordinary $\mathrm{ARMA}(p,q)$ with a seasonal
$\mathrm{ARMA}(P,Q)_s$ gives the \textbf{multiplicative seasonal ARMA}
\[
\Phi(B)\,\Phi^{(s)}(B)\,X_t=\Theta(B)\,\Theta^{(s)}(B)\,\eps_t,
\]
with the four polynomials
\(\Phi(z)=1-\sum_{j=1}^p\phi_jz^j\),
\(\Theta(z)=1-\sum_{k=1}^q\theta_kz^k\),
\(\Phi^{(s)}(z)=1-\sum_{j=1}^P\phi^{(s)}_jz^{sj}\),
\(\Theta^{(s)}(z)=1-\sum_{k=1}^Q\theta^{(s)}_kz^{sk}\).
Special cases: $\mathrm{SMA}(Q)_s$ is $\mathrm{SARMA}(0,0)\times(0,Q)_s$ and
$\mathrm{SAR}(P)_s$ is $\mathrm{SARMA}(0,0)\times(P,0)_s$. Example, a
$\mathrm{SARMA}(0,1)\times(0,1)_{12}$:
\[
X_t=(1-\theta B)\bigl(1-\theta^{(12)}B^{12}\bigr)\eps_t .
\]
\end{definition}

\begin{definition}{ARIMA$(p,d,q)$}{p5-arima}
$\{X_t\}$ is an \textbf{ARIMA$(p,d,q)$} (``integrated ARMA'') process if its
$d$-th difference is a (causal, invertible) $\mathrm{ARMA}(p,q)$:
\[
\Phi(B)\,(1-B)^{d}X_t=\Theta(B)\,\eps_t .
\]
Equivalently, writing $Z_t=\diff^{d}X_t$,
\[
Z_t=\sum_{j=1}^{p}\phi_jZ_{t-j}+\eps_t-\sum_{j=1}^{q}\theta_j\eps_{t-j}.
\]
The ``I'' stands for \emph{integrated}: $X_t$ is recovered from $Z_t$ by summing
(the inverse of differencing) $d$ times. ARMA is the case $d=0$.
\end{definition}

\begin{definition}{SARIMA$(p,d,q)\times(P,D,Q)_s$}{p5-sarima}
The full \textbf{seasonal ARIMA} folds both an ordinary difference of order $d$
and a seasonal difference of order $D$ in front of a SARMA:
\[
\Phi(B)\,\Phi^{(s)}(B)\,(1-B)^{d}(1-B^{s})^{D}\,X_t
=\Theta(B)\,\Theta^{(s)}(B)\,\eps_t .
\]
Here $d$ removes a polynomial trend, $D$ removes a (possibly stochastic)
seasonal pattern, and the four polynomials capture the remaining short-range and
seasonal dependence.
\end{definition}

\begin{definition}{Random walk and the IMA$(1,1)$/ARI$(1,1)$ models}{p5-rw}
The \textbf{random walk} is the AR(1) at the boundary $\phi=1$,
\[
X_t=X_{t-1}+\eps_t ,
\]
which is \emph{not} stationary, but whose first difference $\diff X_t=\eps_t$ is
white noise: it is an ARIMA$(0,1,0)$. Two minimal one-step-integrated models:
\begin{itemize}
\item \textbf{IMA$(1,1)$} $=$ ARIMA$(0,1,1)$: $Z_t=\diff Y_t=\eps_t-\theta\eps_{t-1}$,
      i.e.\ $Y_t=Y_{t-1}+\eps_t-\theta\eps_{t-1}$;
\item \textbf{ARI$(1,1)$} $=$ ARIMA$(1,1,0)$: $Z_t=\diff Y_t=\phi Z_{t-1}+\eps_t$
      with $\abs\phi<1$, i.e.\ $Y_t-Y_{t-1}=\phi(Y_{t-1}-Y_{t-2})+\eps_t$.
\end{itemize}
Neither $\{Y_t\}$ is stationary.
\end{definition}

% ----------------------------------------------------------------------
\subsection{Theorems \& Propositions}

\begin{proposition}{$\diff$ removes a linear trend; $\diff Y$ stays stationary}{p5-lintrend}
Let $X_t=a+bt+Y_t$ with $\{Y_t\}$ stationary. Then
\[
\diff X_t=X_t-X_{t-1}=b+\diff Y_t ,
\]
so the trend $a+bt$ is reduced to the constant $b$, and the random part
$\diff Y_t=Y_t-Y_{t-1}$ is again stationary. Differencing once more removes the
constant too: $\diff^{2}X_t=\diff^{2}Y_t=Y_t-2Y_{t-1}+Y_{t-2}$.\\[2pt]
\textit{Proof idea:} substitute $a+bt+Y_t$ and cancel:
$a+bt-(a+b(t-1))=b$. Linearity of $\diff$ gives $\diff Y_t$; a finite linear
combination of a stationary process is stationary, hence so is $\diff Y_t$.\\
\textit{Why:} the formal justification for first-differencing a series that
``looks linear''; it converts a non-stationary level into a stationary
increment, up to a known constant that the mean estimate then absorbs.
\end{proposition}

\begin{proposition}{$\diff^{q}$ annihilates a degree-$(q-1)$ polynomial trend}{p5-polytrend}
Let $X_t=\mu_t+Y_t$ with $\mu_t$ a polynomial of degree $q-1$ in $t$. Then
$\diff^{q}\mu_t\equiv 0$, so
\[
\diff^{q}X_t=\diff^{q}Y_t=\sum_{k=0}^{q}\binom{q}{k}(-1)^{k}Y_{t-k}.
\]
\textit{Proof idea:} $\diff$ lowers the degree of a polynomial in $t$ by exactly
one (the leading $t^{m}$ becomes $t^{m}-(t-1)^{m}=m\,t^{m-1}+\cdots$); applying
$\diff$ a total of $q$ times to a degree-$(q-1)$ polynomial drives it to the
zero sequence. Linearity then leaves $\diff^{q}Y_t$, expanded by the binomial
theorem, using the induction fact that $\diff^{k}t^{k-1}=0$.\\
\textit{Why:} tells you exactly how many differences $d=q$ are needed to flatten
a polynomial trend of a given degree, which is how one chooses $d$ in
ARIMA$(p,d,q)$.
\end{proposition}

\begin{theorem}{Combined $\diff\diff_{(12)}$ yields a stationary process}{p5-combined}
Let $X_t=a+bt+s_t+Y_t$, where $\{Y_t\}$ is zero-mean stationary with ACVS
$\acvs_\tau$, $a,b$ are constants and $s_{t+12}=s_t$ has period $12$. Then
\[
\diff\diff_{(12)}X_t=\diff\diff_{(12)}Y_t
=Y_t-Y_{t-1}-Y_{t-12}+Y_{t-13}
\]
is weakly stationary, with mean $0$ and
\[
\begin{aligned}
\Cov\bigl(\diff\diff_{(12)}X_t,\ \diff\diff_{(12)}X_{t-\tau}\bigr)
=\;&\acvs_\tau-\acvs_{\tau+1}-\acvs_{\tau+12}+\acvs_{\tau+13}\\
&-\acvs_{\tau-1}+\acvs_{\tau}+\acvs_{\tau+11}-\acvs_{\tau+12}\\
&-\acvs_{\tau-12}+\acvs_{\tau-11}+\acvs_{\tau}-\acvs_{\tau+1}\\
&+\acvs_{\tau-13}-\acvs_{\tau-12}-\acvs_{\tau-1}+\acvs_{\tau} .
\end{aligned}
\]
\textit{Proof idea:} first $\diff_{(12)}X_t=12b+\diff_{(12)}Y_t$ kills $s_t$ and
turns $bt$ into the constant $12b$; then $\diff$ kills the constant, leaving the
fixed linear combination $Y_t-Y_{t-1}-Y_{t-12}+Y_{t-13}$. Its mean is $0$ and
its covariance, expanded by bilinearity, is a function of $\tau$ \emph{only}
(the four-by-four table above).\\
\textit{Why:} the theoretical backing for the standard ``$\diff_{12}$ then
$\diff$'' pre-processing of monthly data (e.g.\ the Mauna Loa $\mathrm{CO}_2$
series); it shows no $t$-dependence survives.
\end{theorem}

\begin{proposition}{Multiplicative-trend variant: $\diff_{(12)}^{2}$}{p5-multtrend}
If instead $X_t=(a+bt)s_t+Y_t$ with $s_{t+12}=s_t$, then
\[
\diff_{(12)}X_t=12b\,s_{t-12}+Y_t-Y_{t-12},
\qquad
\diff_{(12)}^{2}X_t=Y_t-2Y_{t-12}+Y_{t-24},
\]
which has zero mean and a covariance depending only on $\abs\tau$, hence is
stationary.\\[2pt]
\textit{Proof idea:} $\diff_{(12)}$ once turns $(a+bt)s_t$ into $12b\,s_{t-12}$
(using $s_t=s_{t-12}$); a second $\diff_{(12)}$ removes that periodic term as
well, leaving a fixed combination of $Y$'s.\\
\textit{Why:} warns that when trend and season \emph{multiply} (amplitude grows
with level), a \emph{single seasonal} double-difference --- not
$\diff\diff_{(12)}$ --- is what flattens the series.
\end{proposition}

\begin{proposition}{ACVS of the seasonal MA(1)}{p5-smaacvs}
For $X_t=\eps_t-\theta\eps_{t-s}$ with white-noise variance $\sigma_\eps^2$,
\[
\acvs_\tau=\sigma_\eps^2\Bigl[(1+\theta^{2})\delta_\tau-\theta(\delta_{\tau-s}+\delta_{\tau+s})\Bigr]
=\begin{cases}
(1+\theta^{2})\sigma_\eps^2, & \tau=0,\\
-\theta\,\sigma_\eps^2, & \tau=\pm s,\\
0, &\text{otherwise,}
\end{cases}
\]
so the ACF is $\acf_{\pm s}=-\theta/(1+\theta^{2})$ and zero at every other lag.\\[2pt]
\textit{Proof idea:} apply the white-noise filter rule
$\acvs_\tau=\sigma_\eps^2\sum_j a_j a_{j+\abs\tau}$ to the taps
$a_0=1,\ a_s=-\theta$; only $\tau\in\{0,\pm s\}$ pair surviving noise indices.\\
\textit{Why:} a single nonzero ACF spike at the seasonal lag $s$ is the
fingerprint that flags an $\mathrm{SMA}(1)_s$ component in a correlogram.
\end{proposition}

\begin{proposition}{Random walk is non-stationary; its difference is white noise}{p5-rwprop}
For the random walk $X_t=\sum_{j=1}^{t}\eps_j$ (with $X_0=0$),
\[
\E[X_t]=0,\qquad \Var(X_t)=t\,\sigma_\eps^2,\qquad
\Cov(X_t,X_{t+\tau})=t\,\sigma_\eps^2\ \ (\tau\ge0),
\]
all of which depend on $t$, so $\{X_t\}$ is \textbf{not} stationary; but
$\diff X_t=\eps_t$ is white noise.\\[2pt]
\textit{Proof idea:} the variance of a sum of $t$ uncorrelated unit pieces grows
linearly in $t$; the overlap of $X_t$ and $X_{t+\tau}$ is the first $t$ shared
terms.\\
\textit{Why:} the prototype unit-root process --- exactly the ``$\phi=1$
boundary'' that ARIMA differencing is designed to handle.
\end{proposition}

\begin{theorem}{Second-order growth of IMA$(1,1)$; correlation $\to 1$}{p5-imagrowth}
Let $Y_t=Y_{t-1}+\eps_t-\theta\eps_{t-1}$ with $Y_0=0$. Then for $t\ge1$
\[
Y_t=\eps_t+(1-\theta)\sum_{j=1}^{t-1}\eps_j-\theta\eps_0 ,
\]
hence
\[
\Var(Y_t)=\sigma_\eps^2\bigl[\,1+(1-\theta)^{2}(t-1)+\theta^{2}\,\bigr],
\]
and for $\tau>0$
\[
\Cov(Y_t,Y_{t+\tau})=\sigma_\eps^2\bigl[\,1-\theta+(1-\theta)^{2}(t-1)+\theta^{2}\,\bigr].
\]
As $\tau/t\to0$ the correlation $\Corr(Y_t,Y_{t+\tau})\to1$.\\[2pt]
\textit{Proof idea:} unroll the recursion from $Y_0=0$ to get the explicit
linear combination of $\eps_0,\dots,\eps_t$; read off $\Var$ and $\Cov$ from the
shared coefficients (the bulk $(1-\theta)$ terms dominate). Both grow like
$(1-\theta)^2(t-1)\sigma_\eps^2$, so their ratio $\to1$.\\
\textit{Why:} shows quantitatively that an integrated series has
\emph{variance growing linearly in $t$} and \emph{nearby values almost perfectly
correlated} --- the diagnostic signature of needing a difference.
\end{theorem}

\begin{theorem}{Second-order growth of ARI$(1,1)$; correlation near unity}{p5-arigrowth}
Let $Y_t-Y_{t-1}=\phi(Y_{t-1}-Y_{t-2})+\eps_t$ with $\abs\phi<1$ and $Y_0=0$.
Then
\[
Y_t=\sum_{j=1}^{t}\ \sum_{i=0}^{t-j}\phi^{i}\,\eps_j
=\sum_{j=1}^{t}\frac{1-\phi^{\,t-j+1}}{1-\phi}\,\eps_j ,
\]
so that
\[
\begin{aligned}
\Var(Y_t)&=\sigma_\eps^2\sum_{j=1}^{t}\!\left(\frac{1-\phi^{\,t-j+1}}{1-\phi}\right)^{2},\\[2pt]
\Cov(Y_t,Y_{t+\tau})&=\sigma_\eps^2\sum_{j=1}^{t}
\frac{1-\phi^{\,t-j+1}}{1-\phi}\cdot\frac{1-\phi^{\,t+\tau-j+1}}{1-\phi}
\qquad (\tau>0).
\end{aligned}
\]
Again the correlation between nearby times is close to unity.\\[2pt]
\textit{Proof idea:} the increments $Z_t=\diff Y_t$ form a stationary AR(1),
$Z_t=\sum_{i\ge0}\phi^{i}\eps_{t-i}$; integrating ($Y_t=\sum_{m\le t}Z_m$ from
$0$) and summing the inner geometric series gives the closed coefficients
$(1-\phi^{t-j+1})/(1-\phi)$.\\
\textit{Why:} confirms that even when the increments are a well-behaved
stationary AR(1), \emph{integrating} them recreates the random-walk-like
variance growth and near-unit correlation, so the level $\{Y_t\}$ must be
differenced before fitting.
\end{theorem}

% ----------------------------------------------------------------------
% ----------------------------------------------------------------------
